Both formats are plain JSON text files that other tools can read.

\subsection{Layout File}

A layout (\texttt{.json}) holds a complete study. Table~\ref{tab:layoutkeys} lists
its top-level fields.

\begin{table}[H]
  \centering
  \caption{Layout file fields}
  \label{tab:layoutkeys}
  \begin{tabularx}{\textwidth}{l L}
    \toprule
    \textbf{Field} & \textbf{Contents} \\
    \midrule
    \texttt{version} & Format version, currently 1. \\
    \texttt{project} & Project name. \\
    \texttt{pdf} & Path of the background PDF, or \texttt{null}. \\
    \texttt{scene\_rect} & Drawing area: $x$, $y$, width and height, in drawing units (50 per inch). \\
    \texttt{settings} & \texttt{grid\_visible}, \texttt{grid\_opacity}, \texttt{snap}, \texttt{pdf\_visible}. \\
    \texttt{library} & Equipment Library entries: \texttt{symbol}, \texttt{name} and \texttt{properties}. \\
    \texttt{components} & The drawn equipment (below). \\
    \texttt{connections} & Wires, each \texttt{[from\_id, from\_port, to\_id, to\_port]}. \\
    \texttt{zones} & Zone settings (below). \\
    \bottomrule
  \end{tabularx}
\end{table}

\paragraph{Component.} One breaker from the two-zone demonstration:
\begin{lstlisting}
{"id": "breaker_3", "symbol": "breaker", "name": "CB-F1", "display_id": "CB-F1",
 "properties": {"status": "Closed", "symbol_type": "breaker",
   "trip_coil_1_a": 6, "trip_coil_2_a": 5, "close_coil_a": 10,
   "coil_current_basis": "whole_breaker", "mechanism_arrangement": "shared",
   "pole_count": 3, "trip_coil_1_supply": "DC 1", "trip_coil_2_supply": "DC 2",
   "close_coil_supply": "DC 1", "motor_supply": "DC 1"},
 "position": [450.0, 707.0], "scale": 1.0, "rotation": 90.0,
 "ports": {"left": [-32.0, 0.0], "right": [32.0, 0.0]}}
\end{lstlisting}

Currents are stored in amperes under \code{trip_coil_1_a}, \code{trip_coil_2_a},
\code{close_coil_a}, \code{motor_inrush_a} and \code{motor_run_a}. A bus's or surge
arrester's rating is \code{rating_kv}.
\texttt{ports} gives each terminal's position relative to the symbol, before
rotation and scaling. A bus lists its added connection nodes here as well.

\paragraph{Zone.} The demonstration's bus zone, with its first two load rows:
\begin{lstlisting}
{"id": "bus_zone", "name": "BUS-1 Bus Differential", "kind": "Bus",
 "seed": "bus_1", "boundary_cts": ["ct_5", "ct_4", "ct_3"], "confirmed": true,
 "actions": [
   {"component": "breaker_2", "load": "trip_coil_1_a",
    "stage": "Primary trip", "operation": "all_poles"},
   {"component": "breaker_2", "load": "trip_coil_2_a",
    "stage": "Primary trip", "operation": "all_poles"}],
 "topology_signature": "ct-study-v3:43b4...", "configuration_signature": "74b8..."}
\end{lstlisting}

The two signatures are fingerprints of the drawing and of the zone's settings, taken
when the zone was verified. If either no longer matches, the zone must be
reconfirmed. An \emph{Other} load row has a \texttt{component} beginning with
\texttt{other:} and a \texttt{custom\_load} object holding its \texttt{name} and
\texttt{amps}.

\subsection{Netlist File}

A netlist (\texttt{.json}) describes components and connections only, with no zones or
drawing settings. \menu{File > Import Netlist} accepts files like this:

\begin{lstlisting}
{"schema_version": 1,
 "project_name": "Example",
 "components": [
   {"id": "breaker:breaker_1", "display_id": "CB-HV", "type": "breaker",
    "name": "CB-HV", "properties": {"trip_coil_1_a": 10},
    "position": {"x": 300.0, "y": 175.0}, "rotation_deg": 0.0, "scale_factor": 1.0,
    "ports": ["left", "right"],
    "port_positions": {"left": {"x": -32.0, "y": 0.0}, "right": {"x": 32.0, "y": 0.0}}}],
 "connections": [
   {"id": "connection_1",
    "from": {"component": "breaker:breaker_1", "port": "right"},
    "to":   {"component": "ct:ct_1", "port": "left"}}]}
\end{lstlisting}

{\raggedright Each component \texttt{type} is one of the symbol types of
Figure~\ref{fig:symbols}, using the names \code{breaker}, \code{disconnect},
\code{mos}, \code{ct}, \code{pt}, \code{bus}, \code{xfmr_2w}, \code{xfmr_3w},
\code{incoming_line}, \code{surge_arrester}, \code{line_trap}, \code{ground},
\code{junction} and \code{custom}.\par} Every
connection must refer to components and terminals that exist. A file that breaks
these rules is rejected, with a message explaining why.
